% !TeX root = ../教材题库_全部条目_转写工作稿.tex
% Source: XJ-B1-C05-TH-P159-Y1004
% 第5章 / 5.1 任意角和弧度制 / 教材 p.159 / PDF 第166页
% OCR未校对：公式与文字均可能有误；未完成数学TeX转写；禁止直接发给学生
% 原题及相邻原解截图（核对用）：../题目原图/XJ-B1-C05-TH-P159-Y1004-01.webp
% 本文件只在转写完成后作为题干使用；OCR 条目尚不是可靠题干。
\begin{bankitem}{p.159 思考（含紧邻说明）}{第5章 / 5.1 任意角和弧度制 / 教材 p.159 / PDF 第166页}{XJ-B1-C05-TH-P159-Y1004}{OCR未校对：公式与文字均可能有误；未完成数学TeX转写；禁止直接发给学生}
\begin{ocrdraft}
[教材p.159 / PDF第166页]
\par
？思考角度制与弧度制之间可以换算，那么如何换算呢？
\par
一般地，半径为丆的圆周长为2“，所以周角的弧度数是
\par
2丌厂
\par
=2丌，即360。一2丌rad.
\par
欧拉在1748年出版
\par
于是，可以得到角度制和弧度制之间的换算关系：
\par
的《无穷小分析引论》中
\par
明确提出弧度制思想．这
\par
180。一丌rad，
\par
一思想将线段与弧的度量
\par
rad0，01745rad，
\par
统一起来，大大简化了三
\par
180
\par
角函数的微积分公式及
\par
反过来有：1rad\ocrglyph{00AF}
\par
57．30。一57。18'\ensuremath{\cdot}
\par
计算，
\end{ocrdraft}
\end{bankitem}
